%%%%%%%%%%%%%%%%%%%%%%%%%%%%%%%%%%%%%%%%%%%%%%%%%%%%%%%%%%%%%%%%%%%%%%%%%%%%
% 1. FONTS: Times-style text and math font, matching the main document
%    exactly for visual consistency between thesis and slides.
%    Alternatives (swap this block):
%      Libertine:  \usepackage{libertine} + \usepackage[libertine]{newtxmath}
%      XCharter:   \usepackage{xcharter}
%%%%%%%%%%%%%%%%%%%%%%%%%%%%%%%%%%%%%%%%%%%%%%%%%%%%%%%%%%%%%%%%%%%%%%%%%%%%
\usepackage{newtxtext}

%%%%%%%%%%%%%%%%%%%%%%%%%%%%%%%%%%%%%%%%%%%%%%%%%%%%%%%%%%%%%%%%%%%%%%%%%%%%
% 2. MICROTYPOGRAPHY: subtle spacing tweaks for cleaner justified text.
%%%%%%%%%%%%%%%%%%%%%%%%%%%%%%%%%%%%%%%%%%%%%%%%%%%%%%%%%%%%%%%%%%%%%%%%%%%%
\usepackage{microtype}

%%%%%%%%%%%%%%%%%%%%%%%%%%%%%%%%%%%%%%%%%%%%%%%%%%%%%%%%%%%%%%%%%%%%%%%%%%%%
% 3. LANGUAGE & ENCODING: special characters + hyphenation rules.
%    Kept identical to the main document preamble so both compile
%    consistently if you copy content between them.
%%%%%%%%%%%%%%%%%%%%%%%%%%%%%%%%%%%%%%%%%%%%%%%%%%%%%%%%%%%%%%%%%%%%%%%%%%%%
\usepackage[utf8]{inputenc}
\usepackage[T1]{fontenc}
\usepackage[english]{babel}

%%%%%%%%%%%%%%%%%%%%%%%%%%%%%%%%%%%%%%%%%%%%%%%%%%%%%%%%%%%%%%%%%%%%%%%%%%%%
% 4. MATH: equations, symbols, theorem environments, bold vectors.
%    Identical setup to the main document so equations render exactly
%    the same way on slides as in the thesis.
%%%%%%%%%%%%%%%%%%%%%%%%%%%%%%%%%%%%%%%%%%%%%%%%%%%%%%%%%%%%%%%%%%%%%%%%%%%%
\usepackage{amsmath}
\usepackage{amssymb}
\usepackage{amsthm}
\usepackage{mathtools}
\usepackage{bm}
\usepackage{newtxmath} % load after amsmath; swap if using an alt font (see above)
\let\openbox\relax     % avoids a naming clash between amsthm and another package

%%%%%%%%%%%%%%%%%%%%%%%%%%%%%%%%%%%%%%%%%%%%%%%%%%%%%%%%%%%%%%%%%%%%%%%%%%%%
% 5. GRAPHICS, FIGURES & TABLES: same core packages as the main document,
%    minus the ones rarely needed on slides (subcaption, rotating, lscape).
%%%%%%%%%%%%%%%%%%%%%%%%%%%%%%%%%%%%%%%%%%%%%%%%%%%%%%%%%%%%%%%%%%%%%%%%%%%%
\usepackage{float}
\usepackage{caption}
\usepackage{booktabs}
\usepackage{array}
\usepackage{multirow}
\usepackage{adjustbox}

%%%%%%%%%%%%%%%%%%%%%%%%%%%%%%%%%%%%%%%%%%%%%%%%%%%%%%%%%%%%%%%%%%%%%%%%%%%%
% 6. TEXT DECORATION & LISTS
%%%%%%%%%%%%%%%%%%%%%%%%%%%%%%%%%%%%%%%%%%%%%%%%%%%%%%%%%%%%%%%%%%%%%%%%%%%%
\usepackage{enumitem}
\setlist[itemize,1]{label=\usebeamertemplate{itemize item}}
\setlist[itemize,2]{label=\usebeamertemplate{itemize subitem}}
\setlist[itemize,3]{label=\usebeamertemplate{itemize subsubitem}}

\usepackage{bbding}

%%%%%%%%%%%%%%%%%%%%%%%%%%%%%%%%%%%%%%%%%%%%%%%%%%%%%%%%%%%%%%%%%%%%%%%%%%%%
% 7. DIAGRAMS: TikZ for custom diagrams/flowcharts, identical setup to the
%    main document so figures can be copied between thesis and slides
%    without adjustment.
%%%%%%%%%%%%%%%%%%%%%%%%%%%%%%%%%%%%%%%%%%%%%%%%%%%%%%%%%%%%%%%%%%%%%%%%%%%%
\usepackage{tikz}
\usetikzlibrary{arrows.meta,positioning,fit,calc,decorations.pathreplacing}

%%%%%%%%%%%%%%%%%%%%%%%%%%%%%%%%%%%%%%%%%%%%%%%%%%%%%%%%%%%%%%%%%%%%%%%%%%%%
% 8. POSITIONING & COLORS: textpos allows placing elements at absolute
%    coordinates on a slide. xcolor must load before any \definecolor
%    call below, and before tcolorbox.
%%%%%%%%%%%%%%%%%%%%%%%%%%%%%%%%%%%%%%%%%%%%%%%%%%%%%%%%%%%%%%%%%%%%%%%%%%%%
\usepackage[absolute,overlay]{textpos}
\usepackage{xcolor}
\usepackage[most]{tcolorbox}

%%%%%%%%%%%%%%%%%%%%%%%%%%%%%%%%%%%%%%%%%%%%%%%%%%%%%%%%%%%%%%%%%%%%%%%%%%%%
% 9. BIBLIOGRAPHY: same author-year citation setup as the main document,
%    shared via the same file so \citep{}/\citet{} and the citation style
%    stay perfectly consistent between thesis and slides. Edit only in
%    Front Matter/bibliography_setup.tex.
%%%%%%%%%%%%%%%%%%%%%%%%%%%%%%%%%%%%%%%%%%%%%%%%%%%%%%%%%%%%%%%%%%%%%%%%%%%%
%%%%%%%%%%%%%%%%%%%%%%%%%%%%%%%%%%%%%%%%%%%%%%%%%%%%%%%%%%%%%%%%%%%%%%%%%%%%
% BIBLIOGRAPHY: author-year citations (natbib/plainnat = Harvard style,
%     common in economics).
%
%     HOW CITATIONS WORK (mechanism explained):
%     - All sources live in a separate plain-text file, e.g. references.bib,
%       NOT in this .tex file. Each entry has a unique key (e.g. smith2020)
%       plus fields like author, title, year, journal.
%     - \bibliography{references} at the END of your document (without the
%       .bib extension) tells LaTeX which file to pull from.
%     - \bibliographystyle{plainnat} below controls the FORMAT of both the
%       in-text citations and the final reference list. You never type the
%       full reference yourself; it's auto-generated from the .bib fields.
%     - \citep{smith2020}  ->  "(Smith, 2020)"   use when the name is NOT
%       part of your sentence grammar.
%     - \citet{smith2020}  ->  "Smith (2020)"    use when the name flows
%       into the sentence, e.g. "As \citet{smith2020} argues..."
%
%     ABBREVIATING A FREQUENTLY-CITED SOURCE:
%     If you cite the same paper repeatedly (e.g. a key theory paper),
%     define a shortcut once in the preamble:
%       \newcommand{\smithcite}{\citep{smith2020}}
%     Then just type \smithcite throughout instead of retyping the full
%     \citep{smith2020} each time. Useful if you later want to change how
%     that one source is cited everywhere at once.
%
%     COSMETIC APA-LIKE FALLBACK (stay on natbib/bibtex, no engine switch):
%     If a professor asks for "APA-ish" formatting but you don't want to
%     deal with biblatex+biber, just change the style below to:
%       \bibliographystyle{apalike}
%     Everything else (citep/citet, bibtex compilation) stays identical.
%
%     FULL APA 7 / CHICAGO AUTHOR-DATE (more accurate, bigger switch):
%     Requires replacing natbib entirely with biblatex + biber backend and
%     renaming citation commands (citep -> autocite/parencite,
%     citet -> textcite). Only worth it if the professor requires the
%     precise official style rather than something "close enough".
%%%%%%%%%%%%%%%%%%%%%%%%%%%%%%%%%%%%%%%%%%%%%%%%%%%%%%%%%%%%%%%%%%%%%%%%%%%%
\usepackage[round,authoryear]{natbib}
\bibliographystyle{plainnat}  % change to apalike for a quick APA-like look
\setlength{\bibhang}{1em}
\renewcommand{\bibfont}{\small}
\setlength{\bibsep}{1.2em}

%%%%%%%%%%%%%%%%%%%%%%%%%%%%%%%%%%%%%%%%%%%%%%%%%%%%%%%%%%%%%%%%%%%%%%%%%%%%
% 10. BEAMER BEHAVIOUR: hides the default navigation arrows/symbols at the
%     bottom of every slide.
%%%%%%%%%%%%%%%%%%%%%%%%%%%%%%%%%%%%%%%%%%%%%%%%%%%%%%%%%%%%%%%%%%%%%%%%%%%%
\beamertemplatenavigationsymbolsempty

%% =====================================================================
%% 11. UNIVERSITY / PRESENTATION CONFIGURATION
%%  --> Edit ONLY this block to adapt the template to a new university
%%      or department, and to re-brand the theme colour or agenda labels.
%%
%%  NOTE: The presentation title, author, and place are NOT set here.
%%  Edit them directly in Titlepage_Options/Titlepage_slides.tex, inside the
%%  \title{}, \author{}, and \institute{} commands. Add extra lines (e.g. a supervisor's address inside \institute{}) if required.
%% =====================================================================

% --- Institution (reused across all title page options and slide corners) ---
\newcommand{\UniversityName}{Your University}
\newcommand{\UniversityLogo}{images/university_logo.png}
\newcommand{\DepartmentName}{Department / Chair Name}

% --- Agenda / chapter labels used in the progress footline ---
\def\ChapterList{Chapter1,Chapter2,Chapter3,Chapter4,Chapter5,Appendix}

% --- Theme colour (change hex/RGB to re-brand) ---
\definecolor{ThemeColor}{RGB}{43,134,75}

%% =====================================================================
%%  THEME SETUP  (do not need to edit below this line for a new university)
%% =====================================================================

%%%%%%%%%%%%%%%%%%%%%%%%%%%%%%%%%%%%%%%%%%%%%%%%%%%%%%%%%%%%%%%%%%%%%%%%%%%%
% 12. THEME & COLOURS
%%%%%%%%%%%%%%%%%%%%%%%%%%%%%%%%%%%%%%%%%%%%%%%%%%%%%%%%%%%%%%%%%%%%%%%%%%%%
\usetheme{default}
\usecolortheme{Theme}

\setbeamertemplate{section in toc}[sections numbered]

%%%%%%%%%%%%%%%%%%%%%%%%%%%%%%%%%%%%%%%%%%%%%%%%%%%%%%%%%%%%%%%%%%%%%%%%%%%%
% 13. HIGHLIGHTED BOXES
%%%%%%%%%%%%%%%%%%%%%%%%%%%%%%%%%%%%%%%%%%%%%%%%%%%%%%%%%%%%%%%%%%%%%%%%%%%%
\tcbset{
  insightsbox/.style={
    enhanced,
    colback=ThemeColor!4,
    colframe=ThemeColor!85!black,
    boxrule=0.9pt,
    arc=2pt,
    left=6pt,right=6pt,top=6pt,bottom=6pt,
    fonttitle=\bfseries,
    coltitle=white,
    title={Implication},
  }
}

%%%%%%%%%%%%%%%%%%%%%%%%%%%%%%%%%%%%%%%%%%%%%%%%%%%%%%%%%%%%%%%%%%%%%%%%%%%%
% 14. FOOTLINE
%%%%%%%%%%%%%%%%%%%%%%%%%%%%%%%%%%%%%%%%%%%%%%%%%%%%%%%%%%%%%%%%%%%%%%%%%%%%
\usepackage{xstring}
\usepackage{etoolbox}
\newcounter{phasecount}
\newcounter{currentphase}
\setcounter{currentphase}{1}

\newcommand{\ProcessChapter}[1]{%
  \stepcounter{phasecount}%
  \ifnum\thephasecount=\thecurrentphase
    \textcolor{ThemeColor}{\textbf{#1}}%
  \else
    #1%
  \fi
  \ifnum\thephasecount=6\relax\else\ $\rightarrow$ \fi
}

%---- OPTIONAL: Exclude specific slides (e.g. title page, appendix,
%     bibliography) from the frame count shown in the footline ----%
% By default, the footline shows Beamer's built-in \insertframenumber and
% \inserttotalframenumber, which count every single frame in the document,
% including the title page, appendix, and bibliography slides.
%
% If you want an ACCURATE count that only includes your actual content
% slides, uncomment the three lines below AND manually add
% \stepcounter{realframe} inside every content frame that should count
% (i.e. every frame EXCEPT the title page, appendix, and bibliography).
%
% IMPORTANT: after activating this or changing which slides count, you
% need to compile TWICE for the total to display correctly. The first
% compile just records the count; the second one reads it back correctly.
% If the total looks wrong or blank, just recompile once more.
%
% \usepackage{totcount}
% \newtotcounter{realframe}
% \renewcommand{\insertframenumber}{\total{realframe}}
% \newcommand{\realtotal}{\total{realframe}}
% % Then replace \inserttotalframenumber with \realtotal in the footline
% % template below, and add \stepcounter{realframe} inside each counted frame.

\defbeamertemplate{footline}{progressfootline}{%
  \leavevmode\hbox{%
    \begin{beamercolorbox}[wd=.85\paperwidth,ht=2.5ex,dp=1.2ex,left]{}%
      \hspace{1em}\tiny
      \setcounter{phasecount}{0}%
      \edef\temp{\noexpand\forcsvlist{\noexpand\ProcessChapter}{\ChapterList}}%
      \temp
    \end{beamercolorbox}%
    \begin{beamercolorbox}[wd=.14\paperwidth,ht=2.5ex,dp=1.2ex,right]{}%
      \tiny\insertframenumber/\inserttotalframenumber\hspace{0.3em}%
    \end{beamercolorbox}%
  }%
}
\setbeamertemplate{footline}[progressfootline]

\newcommand{\SetFootline}[1]{%
  \setcounter{currentphase}{#1}%
}